\documentclass[12pt,a4paper]{article}

\usepackage[T1]{fontenc}
\usepackage[utf8]{inputenc}
\usepackage{geometry}
\usepackage{setspace}
\usepackage{booktabs}
\usepackage{tabularx}
\usepackage{array}
\usepackage{enumitem}
\usepackage{amsmath}
\usepackage[authoryear,round]{natbib}
\usepackage{hyperref}

\geometry{left=2.8cm,right=2.8cm,top=2.6cm,bottom=2.6cm}
\emergencystretch=2em
\onehalfspacing
\hypersetup{colorlinks=true,linkcolor=black,citecolor=black,urlcolor=black}
\setlist[itemize]{itemsep=0.25em,topsep=0.35em,leftmargin=1.8em}
\setlist[enumerate]{itemsep=0.25em,topsep=0.35em,leftmargin=2.0em}
\newcolumntype{Y}{>{\raggedright\arraybackslash}X}

\title{Procedural Verifiability for Legal AI Outputs:\\
An Audit Framework for Case Recommendation}

\author{}

\date{}

\begin{document}
\maketitle

\begin{abstract}
Legal AI systems can decide which legal authorities become visible before legal reasons are written. This article argues that such visibility power requires a procedural-status test: legal AI outputs should qualify by their verifiable legal-material chain, not by model fluency, autonomy or retrieval accuracy alone. It introduces \textit{normative material screening outputs}: outputs that filter, rank or summarize legal authorities without becoming legal sources or decision reasons. The category is formalized through output effect, system-role caps and a six-part audit vector covering source and corpus, query and version, authority hierarchy, ranking and counter-materials, contestability and adoption logging. The model is implemented in a provider-agnostic harness and evaluated through 230 scenario packets with 609 embedded records or outputs, 51,643 formal invariant checks, 185 metric-separation evaluations, 2,000 metric bootstrap/permutation resamples, 288 counterfactual gate ablations, 4,418 repair-frontier counterfactuals, 355 jurisdiction-profile checks, 820 rank-window visibility checks, 70 rank-order visibility counterfactuals, 57,500 annotation-uncertainty perturbations and 2,990 status-certificate replay checks. The main result is an executable separation result: upstream authority coverage is insufficient for procedural qualification unless the output also preserves a reconstructable source chain, authority set, counter-material pathway and adoption record. Recall above 0.8 predicted procedural qualification with precision 0.25, while the full audit gate cascade achieved precision and specificity of 1.00 on the metric-complete packets. Across all 46 qualified packets, removing evidence, source-tag, authority, counter-material, review, role-cap or adoption gates caused 288/288 counterfactual downgrades. Across 176 blocked high-status claims, all were repairable by declared artefact-gate families. Rank-order visibility counterfactuals downgraded 70/70 packets while preserving authority coverage and counter-material recall. The uncertainty layer produced 0.936 sample-level status stability and 0.928 high-status stability among qualified packets. Thirty endpoint-matched public retrieval outputs remained reference information; ten raw GPT-5.5 outputs with full authority coverage were capped for lack of procedural-source binding; source-supported model-output repairs qualified only after manifest, locator, source-support, rank-salience, contestability and logging gates were satisfied; and 230/230 status certificates replayed the hash, score, role-cap, failure-cap, metric and final-status fields.
\end{abstract}

\noindent\textbf{Keywords:} legal information retrieval; case recommendation; legal AI auditing; contestability; output evidence packets; normative material screening

\section{Introduction}

Legal AI systems increasingly decide what law becomes visible. A case recommender, generated research assistant or court-facing authority report may not decide a dispute, but it can set the field of authorities from which reasons are built. The central risk is not only hallucinated citations or poor retrieval; it is that an output quietly organizes the legal materials available for argument while being treated as neutral search.

This article's claim is direct: procedural status should follow auditability, not intelligence. A fluent answer with weak sources remains reference information. A high-scoring ranked list with no counter-material pathway remains professional support. Only a reconstructable and contestable authority-screening chain can become procedural material. The missing category is \textit{normative material screening}: outputs that filter, rank or summarize legal authorities without becoming legal sources or decision reasons.

The paper therefore asks three operational questions:

\begin{enumerate}[label=RQ\arabic*:]
\item How should AI-generated legal outputs be classified according to their procedural effects?
\item What verifiability and contestability requirements are needed for each output class?
\item How can these requirements be operationalized and stress-tested across implementation, public-output, model-output, source-grounding, adversarial, recoding and uncertainty settings?
\end{enumerate}

The resulting framework has two dimensions. The output-effect dimension classifies what the output does in the legal process. The system-role dimension classifies how the system enters that process. The two dimensions are then linked to verifiability, contestability and accountability requirements. The model is intentionally output-first: it asks when a legal AI output may move from internal reference to procedural material, and when it must be downgraded or withdrawn.

This article contributes a concept, a model and an artifact. The concept is normative material screening: the intermediate procedural role of AI-generated authority lists, ranked case sets and legal-material summaries. The model separates output effect from system role and links both to a six-part audit vector with status thresholds, role caps and failure caps. The artifact is a provider-agnostic harness that implements the model and tests the paper's central contrast: public search outputs can be source-reconstructable but procedurally underqualified; raw model outputs can identify the right authorities yet fail for lack of procedural source binding; source-supported repairs can qualify when source anchors, locators and counter-materials are validated; adversarial repairs are rejected when source, locator, claim, link or counter-material gates fail; blocked outputs receive minimal repair frontiers; and score-uncertainty analysis identifies boundary cases rather than hiding threshold fragility.

AI-based case recommendation is the core case because it sits at the boundary between legal information retrieval and legal reasoning. Case recommendation systems are central to case-based reasoning and legal information retrieval, both long-standing areas of AI and Law research \citep{ashley1992,rissland1995,vanopijnen2017}. They are also procedurally sensitive. In common-law precedent search, civil-law court-decision databases, administrative adjudication systems, arbitral award repositories and ODR platforms, ranking and omission can affect the legal materials available for argument. The framework is therefore jurisdiction-parameterized rather than jurisdiction-neutral: each legal system supplies its own authority hierarchy, while the verifiability problem remains shared.

\section{Related Work and Gap}

\subsection{Legal information retrieval and case-based reasoning}

Legal information retrieval and case-based reasoning both show why case recommendation cannot be evaluated as generic text ranking. Legal materials do not merely describe an external domain; statutes, precedents, administrative rules and court decisions may constitute or structure the legal domain itself. Legal retrieval must therefore account for hierarchy, temporality, authority, citations and professional use rather than reduce relevance to algorithmic similarity \citep{turtle1995,rissland1995,vanopijnen2017}. Case-based reasoning reaches the same point from legal reasoning: cases matter through legally relevant similarities, differences, factors, values and precedential relations \citep{ashley1992,benchcapon2003}. The procedural risk is that a system may make some authorities salient and others invisible while presenting that ordering as ordinary search.

\subsection{Benchmarks and legal retrieval models}

Benchmarking work has substantially improved the technical evaluation of legal retrieval. COLIEE defines shared tasks for case-law retrieval, case-law entailment, statute retrieval and statutory entailment, with explicit datasets, runs and leaderboard metrics \citep{goebel2024}. Neural retrieval systems such as BERT-PLI model paragraph-level interactions between a query case and candidate cases, addressing the length and structure of legal cases more directly than generic ad hoc retrieval models \citep{shao2020}. Legal language and holding benchmarks such as LexGLUE and CaseHOLD extend evaluation to standardized legal-language understanding and legal holding selection \citep{chalkidis2022,zheng2021}. Recent retrieval-augmented legal benchmarks, including reasoning-focused legal retrieval and CLERC, move closer to professional research tasks by connecting retrieval to downstream legal analysis and citation support \citep{zheng2025,hou2025}. LegalBench further extends the evaluation agenda to a broader set of legal reasoning tasks for large language models \citep{guha2023}.

These benchmarks are indispensable, but they answer a different question. COLIEE, LegalBench, CLERC and retrieval benchmarks ask whether a system retrieves, entails, cites or reasons correctly under task conditions. This article asks whether a particular output is procedurally eligible for use after it has been produced. Precision, recall, F1, mean reciprocal rank and entailment accuracy do not determine whether a ranked set may be attached to a court report, relied upon in a professional opinion, or incorporated into a decision reason. Procedural eligibility requires corpus records, authority labels, counter-authority handling, contestability and adoption logs that ordinary retrieval metrics do not measure.

\subsection{Generated legal outputs and source verification}

Generated legal interfaces make the procedural gap more important. A ranked list exposes a result set for inspection; a generated answer may compress selected materials into a fluent narrative and conceal omissions, conflicts or weak source support. Empirical work on legal hallucinations and commercial legal research tools shows that legal language fluency and retrieval augmentation cannot substitute for source verification \citep{dahl2024,magesh2025}. The audit task is therefore to distinguish source truth, retrieval relevance, authority status and procedural use: an output can be technically useful yet unfit for external procedural status, or partly opaque at the model level yet usable if the legal-material chain is reconstructable and contestable.

\subsection{Auditability, explainability and contestability}

The framework also draws on algorithmic accountability and AI auditing literature. Work on accountable algorithms emphasizes the need to review data, models and decision pathways rather than accept automated outputs as self-justifying \citep{kroll2017}. AI auditing scholarship converts this concern into institutional practices for examining systems before and after deployment \citep{raji2020,mokander2023}. Legal and social-science work on explanation and contestability further shows that explanation matters because affected persons need a meaningful opportunity to understand and challenge consequential decisions \citep{wachter2017,binns2018}.

Existing audit and explainability accounts, however, usually address automated decision systems in general. Due-process and public-sector accountability work shows why affected persons need a practical ability to challenge consequential automated support rather than receive a generic explanation after the fact \citep{citron2014,veale2018}. That insight has not yet been translated into a domain-specific account of legal AI outputs that structure the visibility and ranking of legal authorities without themselves deciding a case. This article fills that gap by treating case recommendation as normative material screening and by specifying the audit artefacts needed to qualify such outputs for procedural use.

\section{Case Retrieval as Normative Material Screening}

This article calls such outputs \textit{normative material screening outputs}. The category is not a softer name for retrieval. It identifies an intermediate procedural object: an output that filters, ranks or summarizes legal materials that may later support legal argument or decision-making, without itself creating legal authority or deciding rights and duties. Examples include AI-generated lists of similar cases, recommended statutes, ranked precedent sets, extracted holdings and counter-authority reports.

The distinction also blocks two weak shortcuts. Retrieval metrics such as precision, recall, F1 and mean reciprocal rank are necessary for system development, but they do not say whether a recommendation may be cited in a procedural report, relied upon in a judicial memorandum or incorporated into a decision reason. At the same time, not every legal AI output that influences legal work is a decision-support reason. Many outputs shape the preparation of legal reasoning without becoming reasons for a decision. The middle layer is where most case recommendation systems actually operate.

\section{A Two-Dimensional Output Qualification Model}

The proposed framework starts from a simple premise: the normative qualification of a legal AI output depends on both what the output does and how the system enters the procedure. These are separate dimensions.

\subsection{Output-effect dimension}

The output-effect dimension classifies the legal effect claimed for the output:

The classification uses four criteria: whether the output is exposed outside the user's internal workspace, whether a professional relies on it in advice or preparation, whether it filters the set of legal authorities available for argument, and whether an authorized decision-maker adopts it as part of a reasoned outcome. These criteria keep procedural status separate from model quality.

\begin{enumerate}
\item \textbf{Reference information.} The output supports internal search, summarization or drafting. It has no external procedural effect and is not presented as a basis for another person's legal position.
\item \textbf{Professional support output.} The output assists a lawyer, expert, mediator or other professional in forming advice. The professional remains responsible for verification and communication.
\item \textbf{Normative material screening output.} The output filters, ranks, summarizes or highlights legal authorities or other normative materials. It may shape the materials available for argument, response or judicial consideration.
\item \textbf{Decision-support reason.} The output is adopted, summarized or incorporated by an authorized human or institution as part of the reasons for a legal decision, settlement recommendation, administrative action or arbitral outcome.
\item \textbf{Failure consequence: no external legal effect.} This is not a substantive output class. It is the cap applied when core verifiability, contestability or accountability requirements are absent.
\end{enumerate}

\subsection{System-role dimension}

The system-role dimension determines how far an output may travel procedurally. It is not a second label pasted onto the same output. In the executable protocol, $R(o)$ is a role cap: a back-office tool may produce reference information, a disclosed assistance tool may produce professional support, an auditable procedural tool may produce normative material screening, and an authorized decision-support tool may produce decision-support reasons only after the adoption gate is satisfied. An unaccountable external disposition tool is capped at no external legal effect. The same ranked case list may therefore be private reference, professional support, court-facing screening, or a decision reason depending on the role and adoption path.

\begin{table}[htbp]
\centering
\caption{Output status and procedural role}
\renewcommand{\arraystretch}{1.2}
\begin{tabularx}{\textwidth}{>{\bfseries}p{2.7cm}YYY}
\toprule
Output class & Typical use & Minimum procedural question & Audit artefact required \\
\midrule
Reference information & Internal search, summary, first-pass drafting & Can the user verify the source before relying on it externally? & Source record and user-visible citation or document link. \\
Professional support output & Lawyer advice, expert preparation, mediator support & Can the professional explain the basis and limitations of the output to the affected person? & Source, query, output and limitation record. \\
Normative material screening output & Case recommendation, authority ranking, statute or precedent set selection & Can parties or reviewers reconstruct the corpus, version, ranking logic and omitted counter-materials? & Corpus manifest, ranking record, counter-material record and review log. \\
Decision-support reason & Output adopted in a judgment, administrative decision, arbitral award or ODR recommendation & Can the adoption path, human reasoning and contestation record be audited? & Adoption log, human reasons, contestation record and final responsibility node. \\
No external legal effect & Unaccountable or unverifiable output affecting rights, duties or procedural opportunities & Is any core requirement missing so that use must be blocked or downgraded? & Failure record stating the missing gate. \\
\bottomrule
\end{tabularx}
\end{table}

\section{A Verifiability-Based Audit Model}

The framework can be expressed as an audit model. Let $o$ be a legal AI recommendation output generated in a particular institutional setting. Its procedural status $PS(o)$ is a function of output effect, system role and an audit vector:

\[
PS(o)=g(E(o),R(o),\mathcal{A}(o)).
\]

Here $E(o)$ denotes output effect, $R(o)$ denotes system role, and $\mathcal{A}(o)$ denotes a six-part audit vector:

\[
\mathcal{A}(o)=(S,Q,H,K,T,L).
\]

The six components are source and corpus verifiability ($S$), query and version traceability ($Q$), authority-hierarchy representation ($H$), ranking and counter-material verifiability ($K$), contestability support ($T$), and adoption logging with audit responsibility ($L$). Each component is scored on a three-point ordinal scale: 0 means absent, 1 means partially available, and 2 means procedurally reviewable. The model is not a predictive model of legal correctness. It is a status-allocation protocol: it determines the highest procedural status that the output may claim given the audit evidence available.

\subsection{Formal definitions and invariants}

Let the procedural statuses be ordered as
\[
\mathcal{P}=\{p_0,p_1,p_2,p_3,p_4\},
\]
where $p_0$ is no external legal effect, $p_1$ is reference information, $p_2$ is professional support output, $p_3$ is normative material screening output and $p_4$ is decision-support reason. The order $p_0\prec p_1\prec p_2\prec p_3\prec p_4$ expresses the maximum procedural status an output may claim. Let $\mathcal{A}(o)\in\{0,1,2\}^{6}$ be the audit vector and let $\Pi=(\tau_N,\tau_D,\Gamma,\Delta)$ be a status policy, where $\tau_N$ is the normative-screening threshold, $\tau_D$ is the decision-support threshold, $\Gamma$ is the set of necessary gate predicates and $\Delta$ is the failure-flag cap.

System role is coded as $R(o)\in\mathcal{R}$, where $\mathcal{R}$ contains back-office tools, disclosed assistance tools, auditable procedural tools, authorized decision-support tools and unaccountable external disposition tools. Define a role-cap function
\[
c(R)=
\begin{cases}
p_1,&R=\text{back-office tool},\\
p_2,&R=\text{disclosed assistance tool},\\
p_3,&R=\text{auditable procedural tool},\\
p_4,&R=\text{authorized decision-support tool},\\
p_0,&R=\text{unaccountable external disposition tool}.
\end{cases}
\]
Output effect sets the claimed status; system role sets the maximum status available in the procedure. The operational rule then computes the highest procedurally available status from the audit vector, the role cap and the relevant adoption gate. For a candidate status $p$, define a gate predicate $m_p(\mathcal{A},G_D,\Pi)\in\{0,1\}$. The preliminary status is the greatest status whose audit gate is satisfied:
\[
\sigma_{\Pi}(o)=\max_{\preceq}\{p\in\mathcal{P}:m_p(\mathcal{A}(o),G_D(o),\Pi)=1\}.
\]
Let $G_D=1$ only when the review gate records completed review, authorized adoption, human authorization, jurisdiction assumptions, adoption reasons and contestation record. Failure flags then apply a cap operator $\kappa_F:\mathcal{P}\rightarrow\mathcal{P}$:
\[
PS_{\Pi}(o)=\min_{\preceq}\{c(R(o)),\kappa_{F(o)}(\sigma_{\Pi}(o))\}.
\]
Warning flags leave the status unchanged while recording a defect. Downgrade and suspension flags cap external status at reference information unless the missing artefact is restored. Withdrawal flags cap status at $p_0$.

The protocol has three invariants. First, \textit{gated monotonicity}: if $\mathcal{A}'\geq\mathcal{A}$ componentwise and $R$, $G_D$ and $F$ are unchanged, then $\sigma_{\Pi}(\mathcal{A}')\succeq\sigma_{\Pi}(\mathcal{A})$ only inside the same satisfied gate. Extra score cannot compensate for a missing source, counter-material or adoption gate. Second, \textit{cap dominance}: for any output $o$, $PS_{\Pi}(o)\preceq\sigma_{\Pi}(o)$ and $PS_{\Pi}(o)\preceq c(R(o))$, while withdrawal flags force $PS_{\Pi}(o)=p_0$ regardless of upstream recall. Third, \textit{metric non-equivalence}: retrieval metrics may inform $K$ or issue-packet construction, but they are evidence variables, not status variables. These invariants are what make the model non-compensatory: legal fluency, high recall or strong ranking cannot purchase procedural status when a required audit artefact is absent.

Three consequences follow immediately. \textit{Separation}: for any output with perfect upstream authority coverage but a non-procedural source tag, $PS_{\Pi}(o)\preceq p_1$ whenever it claims external screening status, because the source-attribution failure cap applies after $g_0$; missing output-source links are treated the same way. \textit{Screening necessity}: an output cannot reach $p_3$ unless each audit dimension is at least partially present, the output evidence packet is structurally present, and the counter-material gate is declared and reviewable; a high total score alone is insufficient. \textit{Adoption necessity}: an output cannot reach $p_4$ unless $T=L=2$ and $G_D=1$, even if it already qualifies as screening. These are protocol properties, not empirical assumptions; the experiments below test whether the implementation enforces them on public, model-generated, repaired and adversarial packets.

\subsection{Audit dimensions}

\begin{table}[htbp]
\centering
\caption{Audit dimensions for case recommendation outputs}
\renewcommand{\arraystretch}{1.15}
\begin{tabularx}{\textwidth}{>{\bfseries}p{2.5cm}p{2.3cm}YY}
\toprule
Dimension & Score 0 & Score 1 & Score 2 \\
\midrule
$S$: Source and corpus & Sources or corpus cannot be reconstructed. & Individual source links exist, but corpus scope or update status is incomplete. & Source links, corpus manifest, inclusion and exclusion rules, update date and document versions are recorded. \\
$Q$: Query and version & Query, filters or output version are unavailable. & Some query or version information is retained. & Query, filters, issue labels, source-collection version, workflow or interface version and output identifier are reconstructable. \\
$H$: Authority hierarchy & Output does not distinguish legal status. & Output marks broad categories such as binding or persuasive. & Output records jurisdiction, court level, legal force, validity, later treatment and relation to governing norms. \\
$K$: Ranking and counter-materials & Ranking is opaque and contrary materials are not surfaced. & Broad ranking factors or some contrary materials are shown. & Ranking/re-ranking factors, top-$k$ set, omitted-material checks and counter-authority recall are reviewable. \\
$T$: Contestability & Affected persons cannot inspect or challenge the output. & A professional or internal reviewer can challenge the output. & Parties or reviewers can inspect, challenge, supplement and obtain a recorded response. \\
$L$: Adoption and responsibility & No adoption path or responsible actor is recorded. & Output and user action are logged. & Output ID, query, result list, adopted and rejected items, reasons, timestamp and responsible reviewer are logged. \\
\bottomrule
\end{tabularx}
\end{table}

This scoring scheme deliberately avoids requiring full model-parameter transparency in every case. A legal procedure may not need source code to decide whether a ranked case list can be contested. It does need enough information to reconstruct the legal-material chain and to test whether a relevant authority was omitted, misranked, outdated or mischaracterized.

\subsection{Status thresholds}

The audit vector determines the maximum status available to an output. A high aggregate score does not compensate for failure of a necessary gate. The model therefore uses both minimum gates and total-score bands.

\begin{table}[htbp]
\centering
\caption{Status allocation rules}
\renewcommand{\arraystretch}{1.15}
\begin{tabularx}{\textwidth}{>{\bfseries}p{3cm}YY}
\toprule
Candidate status & Necessary gates & Additional condition \\
\midrule
Reference information & $S\geq1$ and $Q\geq1$ & Output remains internal or user verifies sources before external reliance. \\
Professional support output & $S,Q,L\geq1$ & Professional user accepts verification responsibility and can explain limits. \\
Normative material screening output & $S,Q,H,K,T,L\geq1$ & Total score $\sum \mathcal{A}(o)\geq9$ and counter-material handling is at least partially reviewable; declared complete-set omissions trigger caps. \\
Decision-support reason & $S,Q,H,K,T,L\geq1$, $T=L=2$ and authorized-adoption gate satisfied & Total score $\sum \mathcal{A}(o)\geq10$ plus completed review, human authorization, reasons and recorded contestation. \\
No external legal effect & Any necessary gate for the claimed status is missing & Output is blocked, downgraded or retained only as internal reference. \\
\bottomrule
\end{tabularx}
\end{table}

The mapping function $g$ is therefore explicit rather than discretionary. First, assign the highest status whose gates are satisfied:

\[
g_0(\mathcal{A})=
\begin{cases}
p_4, & S,Q,H,K,T,L\geq1,\ \sum\mathcal{A}\geq10,\ T=L=2,\ G_D(o)=1,\\
p_3, & S,Q,H,K,T,L\geq1 \text{ and } \sum \mathcal{A}\geq9,\\
p_2, & S,Q,L\geq1,\\
p_1, & S,Q\geq1,\\
p_0, & \text{otherwise.}
\end{cases}
\]

Second, apply system-role and failure caps. A disclosed assistance tool cannot exceed professional support even if the audit vector would otherwise satisfy normative screening; an auditable procedural tool cannot become a decision-support reason without authorized adoption. Withdrawal-level failures, such as unreconstructable sources, materially false generated summaries or irreversible external action without human authorization, set $PS(o)$ to no external legal effect. Suspension or downgrade flags, such as high-authority omission, invalid-authority recommendation, counter-material suppression, missing source-attribution tags, absent jurisdiction assumptions, review-gate failure, ranking drift or contestation failure, cap the output at reference status until the missing artefact is restored. Thus $g$ combines the gate function $g_0$ with role and failure caps. If $K=0$, the output cannot be treated as normative material screening because the ranking and counter-material treatment cannot be reviewed. If $T<2$, $L<2$ or $G_D(o)=0$, the output cannot be a decision-support reason, even if retrieval accuracy is high. If the system purports to trigger an external disposition without legal authorization, review, contestability or accountability, the output has no external legal effect.

Operationally, an evaluator applies the model in six steps: identify the candidate procedural status claimed for the output; freeze the relevant issue, query, source collection, workflow version and output identifier; score each component of $\mathcal{A}(o)$ using the audit artefacts in Table 2; apply the necessary gates in Table 3; assign the highest available status or record the missing gate; and repeat the audit after material corpus, workflow or interface changes. This makes the framework usable as a deployment checklist, post-use audit and benchmark extension.

The thresholds use non-substitutable gates. The 0--2 scale is ordinal: 0 means absent, 1 means incomplete or internal, and 2 means reviewable by the relevant professional or procedural actor. The default $\tau_N=9$ permits one limited weakness while still requiring every gate; decision-support status requires $T=L=2$ because external reasons require both contestation and adoption records. The full validation suite re-evaluates all 230 packets under thresholds 8--12. Thresholds 8--10 produced identical status allocations; only stricter thresholds generated demotions, and none generated promotions.

\begin{table}[htbp]
\centering
\caption{Full-packet sensitivity of status allocation to threshold policy}
\scriptsize
\renewcommand{\arraystretch}{1.1}
\begin{tabularx}{\textwidth}{rrrrY}
\toprule
$\tau_N$ & $\tau_D$ & Flips & Demotions & Status distribution \\
\midrule
8 & 10 & 0 & 0 & decision 12; screening 34; professional 23; reference 129; no effect 32 \\
9 & 10 & 0 & 0 & decision 12; screening 34; professional 23; reference 129; no effect 32 \\
10 & 11 & 0 & 0 & decision 12; screening 34; professional 23; reference 129; no effect 32 \\
11 & 12 & 2 & 2 & decision 12; screening 32; professional 25; reference 129; no effect 32 \\
12 & 13 & 41 & 41 & decision 0; screening 17; professional 52; reference 129; no effect 32 \\
\bottomrule
\end{tabularx}
\end{table}

\subsection{Jurisdictional parameterization}

The audit dimensions are stable, but the materials satisfying $H$ and $K$ are local. For each setting $j$, evaluators instantiate authority hierarchy as $H_j$ and counter-material handling as $K_j$ before scoring. A high-status output must also carry jurisdiction assumptions that support the declared profile; an otherwise complete packet is downgraded when the profile is absent or mismatched.

\begin{table}[htbp]
\centering
\caption{Jurisdictional profiles for $H_j$ and $K_j$}
\renewcommand{\arraystretch}{1.15}
\begin{tabularx}{\textwidth}{>{\bfseries}p{2.6cm}YY}
\toprule
Setting & $H_j$: authority hierarchy & $K_j$: counter-materials and limiting materials \\
\midrule
Common law & Binding precedent, persuasive precedent, statutes, subsequent treatment, overruled or superseded authorities. & Distinguishing cases, contrary lines, limiting subsequent treatment, minority lines and overruled authorities. \\
Civil law & Constitutional norms, code provisions, special statutes, regulations, high-court decisions, lower-court decisions and doctrinal commentary. & Competing statutory interpretations, systematic conflicts, teleological limitations, contrary doctrinal views and later amendments. \\
Admin. adjudication & Statutes, binding regulations, agency rules, official guidance, adjudicative decisions and non-binding policy documents. & Conflicting agency interpretations, withdrawn guidance, higher-level statutory constraints and contrary adjudicative decisions. \\
Arbitration & Contract text, chosen law, mandatory law, institutional rules, prior awards and trade usage. & Contrary awards, mandatory-law limits, distinguishing clauses, procedural fairness objections and public-policy exceptions. \\
ODR & Platform rules, consumer statutes, terms of service, regulatory guidance and prior platform decisions. & Consumer-protection overrides, exception policies, appeal decisions and human-review overrides. \\
\bottomrule
\end{tabularx}
\end{table}

\subsection{Contestability and audit responsibility}

Contestability requires a practical opportunity to inspect selected materials, identify missing or contrary materials, submit alternatives, request human review and receive a recorded response when the output is adopted. Adoption logging is distinct: it records whether, why and by whom the output was taken up after contestation was possible. Audit responsibility assigns each artefact to the actor that controls it: developers for capability and safety claims, deployers for corpus and interface records, professional users for issue framing and verification, and authorized adopters for decision uptake. The point is status evidence, not a general liability regime.

\section{Operationalizing the Framework for Case Recommendation}

AI-based case recommendation can be produced by many upstream architectures, but the audit framework does not depend on their internal mechanics. The legally relevant object is the output evidence packet: the source collection, versions, legal materials, output units, source locators, authority labels, counter-material links and adoption records that make the output reconstructable. Failures in these artefacts create different legal consequences. A stale source collection creates source and temporal risk. An ordering rule that privileges surface similarity may miss legally decisive distinctions. A generated summary may misstate holdings or flatten procedural posture. A user interface may display confirming cases while hiding contrary authorities.

The evaluation protocol should therefore separate upstream performance metrics from procedural checks. Precision, recall, F1, MRR and related metrics may help evaluate an upstream retrieval or recommendation component. They do not evaluate whether the output has preserved the legal hierarchy of authorities, displayed counter-materials, distinguished binding and persuasive materials, recorded update status, or produced a contestable adoption trail. These are evaluation requirements for procedural status.

The procedural metrics can be stated directly. \textit{Authority coverage} asks whether known high-authority materials for the issue appear in the output set. \textit{Counter-authority recall} asks whether legally relevant contrary or limiting materials are surfaced. \textit{Hierarchy accuracy} asks whether the output correctly labels binding, precedential, persuasive, overruled or superseded materials. \textit{Version traceability} asks whether the source collection and upstream run state that produced the output can be reconstructed. \textit{Procedural-source coverage} asks whether source links are bound to official, tool-verified or otherwise procedurally usable source tags rather than merely named. \textit{Adoption-log completeness} asks whether later human use of the output is observable. These metrics do not replace upstream benchmarks; they determine whether a benchmarked output can be used procedurally.

\begin{table}[htbp]
\centering
\caption{Decision matrix for legal AI output qualification}
\renewcommand{\arraystretch}{1.2}
\begin{tabularx}{\textwidth}{>{\bfseries}p{3.1cm}YYY}
\toprule
Output use & Minimum verifiability & Minimum contestability & Allowed procedural status \\
\midrule
Internal legal search & Source and version record; user can inspect retrieved materials & Internal user review & Reference information \\
Professional advice support & Source, query, output and limitations recorded & Client or supervising professional can request explanation when output materially shaped advice & Professional support output \\
Court or tribunal case recommendation & Corpus scope, update date, authority hierarchy, ranking factors, counter-authority handling and version record & Parties can inspect, challenge and supplement recommended materials & Normative material screening output \\
Decision reasoning support & Full adoption path, human confirmation, contrary-material treatment and audit trail & Parties can contest use; decision-maker gives reasons for adoption or rejection & Decision-support reason \\
Unaccountable external disposition & Missing legal authorization, review, contestability or accountability chain & Not satisfied & No external legal effect \\
\bottomrule
\end{tabularx}
\end{table}

\subsection{Baseline testing}

Baseline testing should cover both technical and legal dimensions. A case recommendation system should be tested against known authority sets, especially high-authority or binding materials. Missing a binding, precedential, designated or otherwise high-authority material is different from missing one marginally similar low-level case. The former may trigger withdrawal from procedural use even if aggregate retrieval metrics remain acceptable.

Baseline testing should also record database coverage, update intervals, jurisdictional boundaries, language coverage, case hierarchy, citation treatment and exclusion rules. A system trained or indexed on a narrow corpus should not be evaluated as if it were a general legal research system. A system that excludes unpublished decisions, local court decisions or outdated but historically relevant authorities should make that exclusion visible to the user.

\subsection{Counter-material presentation}

Case recommendation is procedurally dangerous when it confirms one line of reasoning without exposing plausible counter-materials. A useful legal recommendation system should therefore include counter-authority handling as a first-order evaluation requirement. The goal is not to force neutrality in the abstract. The goal is to preserve the adversarial and reason-giving structure of legal reasoning by ensuring that contrary cases, limiting cases, overruled authorities and minority lines are visible when they are legally relevant.

For audit purposes, counter-authority recall can be measured against a manually curated issue set:

\[
CAR=\frac{|\text{retrieved counter-authorities}\cap\text{known counter-authorities}|}{|\text{known counter-authorities}|}.
\]

The denominator is not the universe of all contrary law. It is the set of contrary or limiting materials identified for a defined issue by the relevant court, professional reviewer or benchmark curator. This makes the measure imperfect but usable.

This requirement is particularly important for generated legal outputs. A ranked list may visibly display multiple authorities. A generated summary may compress them into a single narrative and thereby hide disagreement. Evaluation should test whether the output cites supporting materials accurately, distinguishes holding from dicta, distinguishes binding from persuasive authorities, and signals unresolved conflict instead of hallucinating consensus. Empirical work on legal hallucinations shows why source verification cannot be left to model fluency \citep{dahl2024}.

\subsection{Output evidence packets}

The audit layer proposed here sits above any particular search, retrieval, database, generation, agent or manual review architecture. The framework therefore does not inspect upstream implementation details. It requires an \textit{output evidence packet}: a provider-agnostic record that connects the legal propositions or recommended materials in an output to source identifiers, versions, locators and review status.

A minimally reviewable evidence packet should record the issue presented to the system, source collection identifier, source version, output identifier, output units, source identifiers linked to each unit, paragraph or section locators, source-attribution tag, authority labels, counter-material links, human review gate, human adoption record and objection record. This requirement is related to model and dataset documentation work, but the unit of documentation is a procedurally relevant legal output rather than a model card or dataset sheet \citep{gebru2021,mitchell2019}. The upstream system may be a keyword search engine, legal database, generated-output workflow, agent workflow or manual review process. The audit question is the same: can the legal-output chain be reconstructed, source-tagged, reviewed and contested?

Four structural packet metrics make this requirement operational. \textit{Structural evidence fidelity} is the proportion of output-source links marked as supporting the output unit to which they are attached:

\[
EF=\frac{|\text{output-source links that support the unit}|}{|\text{all output-source links}|}.
\]

\textit{Evidence coverage} is the proportion of legally material output units that have at least one reconstructable source locator:

\[
EC=\frac{|\text{output units with source-locator support}|}{|\text{all legally material output units}|}.
\]

\textit{Source-tag coverage} measures output-source links carrying a verification-status tag:

\[
STC=\frac{|\text{output-source links with source-attribution tags}|}{|\text{all output-source links}|}.
\]

\textit{Procedural-source coverage} is stricter:

\[
PSC=\frac{|\text{links tagged official, tool-verified, user-verified or settled}|}{|\text{all output-source links}|}.
\]

The tag may indicate whether the source was verified by a connected legal research tool, drawn from public metadata, supplied by the user, treated as settled background, or still requires pinpoint verification. These are packet-integrity metrics, not automated merits determinations: they test whether the output-source chain is structurally reviewable, while doctrinal support still requires professional or institutional review. For external normative screening, a tag that merely says ``needs verification'' records work still to be done. Low $EF$, $EC$, $STC$ or $PSC$ means that the output is not procedurally reconstructable even if an upstream system found relevant documents.

\subsection{Review and adoption gates}

Source reconstruction does not by itself authorize legal reliance. The gate record should state whether professional review is required and completed, what jurisdiction assumptions were used, which reliance level applies, whether any irreversible external action has human authorization, and what the human actor did with the output. An accurate draft remains below decision-support status until an authorized actor reviews it, adopts it and records reasons.

The minimum adoption-log schema is: output identifier, time, user role, issue, filters, source-collection version, workflow version, displayed material set, displayed counter-materials, source-attribution tags, review-gate status, adopted item, rejected item, human reason, objection or supplement submitted, response to objection and final responsibility node. A system that cannot export these fields should remain below decision-support status.

\section{Downgrade and Withdrawal Conditions}

Procedural status is dynamic. Corpus changes, workflow updates, ranking drift, source-attribution gaps, review-gate failures or contestation failures can trigger failure caps after an output or system was initially approved. Table 5 operationalizes those caps: warnings correct isolated defects, downgrades cap status until a missing artefact is restored, suspensions block external use for an issue class, and withdrawals remove external use when the legal-material chain cannot be reconstructed or an irreversible action lacks human authorization.

\begin{table}[htbp]
\centering
\caption{Downgrade severity levels}
\renewcommand{\arraystretch}{1.15}
\begin{tabularx}{\textwidth}{>{\bfseries}p{2.8cm}YY}
\toprule
Severity & Trigger & Procedural consequence \\
\midrule
Warning & Isolated error with reconstructable source, query and adoption record. & Correct the output and monitor the issue class. \\
Downgrade & Missing authority hierarchy, incomplete counter-material display, missing source tags, incomplete review gate or weak adoption logging. & Limit output to reference or professional support status until the missing artefact is restored. \\
Suspension & Repeated omission of high-authority materials, invalid authority recommendation or unexplained ranking drift in an issue class. & Suspend external procedural use for that domain or data source. \\
Withdrawal & Source, corpus, ranking or adoption chain cannot be reconstructed, the system generates materially false legal summaries, or an irreversible action lacks human authorization. & Withdraw from external legal use; retain only internal testing or research use. \\
\bottomrule
\end{tabularx}
\end{table}

\section{Implementation and Evaluation}

To test whether the framework is operational rather than merely classificatory, the rules above were implemented in a lightweight command-line audit harness. The implementation uses local JSON scenario files as the unit of evidence and produces both Markdown and JSON run artefacts. The companion artifact supplies the harness, scenario files, source snapshots, run artefacts and replication commands.

The harness implements the mapping in Section 4 directly. Each scenario declares a claimed status, system role or role-inference evidence, a base audit vector $S,Q,H,K,T,L$, a jurisdiction profile, upstream-performance metrics and expected outcomes. Authority sets, evidence packets, review gates and counter-material declarations are mandatory for outputs that claim or qualify for screening or decision-support status; missing fields trigger caps rather than silent upgrades. The runner does not read the expected outcome when computing status. It computes the allowed status, role cap, counter-authority recall, authority coverage, invalid-authority rate, structural evidence fidelity, evidence coverage, source-tag coverage, procedural-source-tag coverage and downgrade or withdrawal disposition. Separate commands vary status thresholds and re-score the same packets under strict and lenient coding policies.

The complete evaluation matrix contains 230 primary scenario packets with 609 embedded records or outputs, plus 30 public source-text anchor checks, 50 raw model-output transcript locator checks, 51,643 formal invariant checks, 185 metric-separation evaluations, 2,000 metric bootstrap/permutation resamples, 288 gate-ablation evaluations, 4,418 repair-frontier counterfactuals, 355 jurisdiction-profile checks, 820 rank-window visibility checks, 70 rank-order visibility counterfactuals and 2,990 status-certificate replay checks. The evidence is deliberately separated into ten groups. Implementation and invariant tests check whether gates and caps execute as specified. Public-output audits test whether real legal search and listing outputs can be reconstructed and status-capped without model access. Model-output interventions test whether model-output variants change status when source binding, counter-material salience, contestability, logging and adoption gates change. Metric-separation tests compare upstream precision, recall and F1 thresholds with procedural status and resample the key estimates. Gate-ablation tests remove procedural conditions from every qualified packet. Repair-frontier tests compute minimal upgrade paths for blocked high-status claims. Jurisdiction-profile tests mutate assumptions to separate valid cross-profile parameterization from profile mismatch. Ranking-visibility tests compute a window-sensitivity curve for counter-material salience and, where packet structure permits, hold authority coverage and counter-material recall fixed while moving counter-material below the visibility window. Status-certificate tests replay the derivation fields for every allocation. Robustness checks test threshold and coding fragility. The central empirical question is not whether every upstream system is accurate; it is whether the protocol separates retrieval success, source reconstructability and procedural qualification.

\begin{table}[htbp]
\centering
\caption{Claim-to-evidence map}
\small
\renewcommand{\arraystretch}{1.1}
\begin{tabularx}{\textwidth}{>{\bfseries}p{3.2cm}YY}
\toprule
Claim & Evidence layer & What would falsify the claim \\
\midrule
The protocol is executable & Stress scenarios, issue ablations and unit tests & Gates, caps or downgrade rules produce the wrong status. \\
Public outputs are auditable without model access & Metadata, public-system outputs and endpoint-matched retrieval & Visible records cannot be frozen, source-tagged or reconstructed. \\
Retrieval success is not procedural qualification & Public retrieval and raw GPT-5.5 outputs & High authority coverage alone upgrades the output to screening. \\
Retrieval metrics are not procedural status & Metric separation analysis & Precision, recall or F1 thresholds predict screening qualification with high precision. \\
Source-supported repair can upgrade an output & Repaired model outputs, evidence ladder and mixed-authority issue packets & Validated source, locator, counter-material and tag evidence still fails to qualify. \\
The validator rejects brittle repairs & Sixty adversarial repair variants & Locator, source, claim, link or counter-material attacks pass as screening. \\
Manifest support is externally anchorable & Public source-text anchor verification & Manifest support terms cannot be located in public source snapshots. \\
Raw model-output coding is transcript-anchored & Model-output transcript verification & Scenario locators cannot be found in the committed raw transcript. \\
Formal status properties are executable & Formal invariant verification & Scores alone can buy screening status without authority-set, evidence-packet, counter-material, review-gate, role-cap or adoption conditions. \\
Qualified outputs depend on their gates & Gate ablation analysis & Removing procedural gates from qualified packets leaves status unchanged. \\
Blocked outputs have repair frontiers & Repair-frontier analysis & Blocked high-status claims cannot be mapped to minimal source, authority, contestability, role or adoption repairs. \\
Rank salience is independently status-relevant & Ranking-visibility diagnostics & Coverage-preserving rank-order drift leaves high procedural status unchanged. \\
Status allocation is replayable & Status-certificate replay & Scenario-hash, score, role-cap, failure-cap, metric or final-status fields cannot be reproduced. \\
\bottomrule
\end{tabularx}
\end{table}

Derived robustness checks then vary thresholds, recode audit vectors under strict and lenient policies and run score-blinded codebook coding. These checks are not counted as independent field observations; they test whether the same packet evidence leads to stable status allocation.

\begin{table}[htbp]
\centering
\caption{Full evaluation matrix}
\small
\renewcommand{\arraystretch}{1.1}
\begin{tabularx}{\textwidth}{>{\bfseries}p{2.8cm}p{1.5cm}p{2.4cm}Y}
\toprule
Layer & Units & Rule pass & What it checks \\
\midrule
Stress scenarios & 10 & 10/10 & Status gates, failure caps, downgrade and withdrawal behavior. \\
Public metadata & 120 records in 6 files & 6/6 & Source reconstruction across six public legal-record sources. \\
Public-system outputs & 60 records in 6 files & 6/6 & Real upstream legal-output reconstruction and professional-support caps. \\
Issue-specific public output/source packets & 19 records in 3 files & 3/3 & Public issue-search outputs and one source-bound public-source packet tested against high-authority and counter-material requirements. \\
Endpoint-matched public retrieval benchmark & 225 records in 30 files & 30/30 & Public case-law or known-item outputs tested against authority sets that the endpoint can return; mixed authority gaps are recorded separately. \\
Raw model outputs & 10 & 10/10 & High authority coverage without source binding remains procedurally capped. \\
Source-supported model-output repairs & 10 & 10/10 & Same model outputs qualify only after a validator checks manifest membership, locators, hashed source-support excerpts, source tags and issue-set coverage. \\
Model-output evidence ladder & 70 & 70/70 & Same model outputs move across reference, professional-support, screening, decision-support and no-effect status only when evidence and adoption gates change. \\
Adversarial source-support repairs & 60 & 60/60 & Negative controls test locator, claim, contradiction, source, link and counter-material attacks against the repair validator. \\
Public source-text anchors & 30 checks & 30/30 & Support terms are checked against extracted public source text snapshots, reducing manifest-only validation. \\
Model-output transcript anchors & 50 checks & 50/50 & Raw model-output scenario locators are checked against the committed transcript sections. \\
Formal invariant verification & 51,643 checks & 51,643/51,643 & Exhaustive generated states test monotonicity, authority-set necessity, evidence-packet necessity, counter-material necessity, adoption gates, role caps, failure caps and metric non-equivalence. \\
Metric separation & 185 evaluations & Recall-threshold precision 0.25; full gate precision 1.00 & Tests whether upstream precision, recall and F1 predict procedural qualification. \\
Metric statistical robustness & 2,000 resamples & Bootstrap CI for recall-threshold precision 0.19--0.32; permutation $p=0.538$ & Tests whether the retrieval/status separation is stable under resampling and label permutation. \\
Qualified-output gate ablations & 288 ablations & 288/288 & Removes procedural gates from all qualified packets and checks that status falls below the corresponding level. \\
Blocked-claim repair frontiers & 4,418 counterfactuals & 176/176 & Computes minimal artifact-gate families needed to upgrade blocked normative-screening or decision-support claims. \\
Jurisdiction-profile mutations & 355 checks & 138/138 & Valid generic profile assumptions preserve status, while missing or mismatched profile assumptions downgrade qualified packets. \\
Ranking-visibility diagnostics & 820 window checks / 70 counterfactuals & 70/70 & Computes counter-material visibility across rank windows; rank-order counterfactuals preserve authority coverage and counter-material recall while activating ranking-drift caps. \\
Status certificate replay & 2,990 checks & 2,990/2,990 & Replays machine-readable status certificates for all 230 packets from scenario hash, score candidate, role cap, failure cap, metrics and final status. \\
Mixed-authority public source-screening packets & 5 & 5/5 & Source-bound statute, case and source packets can qualify as normative screening. \\
Issue ablations & 20 & 20/20 & High-authority omissions, counter-material suppression, unverified source tags and missing adoption gates trigger caps. \\
Annotation robustness & 460 recodings & 218/230 stable & Strict and lenient recoding policies test status stability across all packets. \\
Annotation uncertainty & 57,500 perturbations & 0.936 sample stability; 0.928 qualified high-status stability & Fixed-seed score-noise perturbations locate threshold-boundary cases and test coding-noise robustness. \\
Score-blinded dual coding & 230 packets, 2 passes & 0.99 coder agreement; 0.95 min base agreement & Coding passes score packets without original scores or expected outcomes. \\
Threshold sensitivity & 1150 evaluations & 0 flips at thresholds 8--10 & All 230 packets are re-evaluated under normative thresholds 8--12. \\
\bottomrule
\end{tabularx}
\end{table}

\begin{table}[htbp]
\centering
\caption{Scenario-based harness results}
\small
\renewcommand{\arraystretch}{1.1}
\begin{tabularx}{\textwidth}{>{\bfseries}p{3.2cm}p{2.1cm}p{2.5cm}p{1.2cm}p{1.2cm}Y}
\toprule
Scenario type & Profile & Allowed status & Score & Up. recall & Audit result \\
\midrule
Court authority report & Common law & Normative screening & 10 & 0.83 & Procedurally usable despite incomplete counter-material recall. \\
Civil-law statutory interpretation & Civil law & Normative screening & 10 & 0.82 & Uses competing interpretations and amendments as counter-materials. \\
Administrative guidance conflict & Admin. & Reference only & 8 & 0.88 & Strong upstream recall but suspended because counter-material and invalid-guidance flags fire. \\
Grounded output summary & Common law & Normative screening & 11 & 0.90 & Output-source mapping and structural evidence fidelity are complete. \\
High-coverage uncontestable output & Common law & No external effect & 5 & 0.93 & Strong upstream coverage does not cure missing contestability and weak evidence links. \\
Authorized ODR review & ODR & Decision-support reason & 12 & 0.91 & Permitted only because authorization, appeal, contestation and adoption logs are present. \\
\bottomrule
\end{tabularx}
\end{table}

Across ten stress-test scenarios, all rule outcomes were reproduced. The mean audit score was 8.30 and mean upstream recall was 0.84. Three scenarios had upstream recall above 0.80 but were procedurally blocked or downgraded. A high-recall output may still be unfit for external legal use when it omits a high-authority material, suppresses a relevant limiting material, recommends invalid authority, lacks output-source evidence or prevents affected persons from contesting the output.

\subsection{Public legal-record reconstruction}

The harness was also run on public legal-record snapshots. One layer sampled 20 public metadata records from each of six legal systems. A second layer froze the top ten records visible in six public legal retrieval or listing streams. These layers test reconstruction: whether visible records can be frozen, source-tagged, linked to locators and assigned a procedural cap by the audit layer. In both layers, all six scenario files were classified as professional support outputs because public record visibility establishes source reconstructability, not issue-specific ranking, counter-material completeness or party-facing contestation.

\subsection{Issue-specific public output/source audit}

A stricter public layer then froze two issue-specific public search outputs and one source-bound public-source packet, converting their visible records into evidence packets. The U.S. packet used a CourtListener issue search for agency deference after \textit{Loper Bright}; the English-law packet used a National Archives issue search for mesothelioma causation after \textit{Fairchild}; and the EU packet used selected Legislation.gov.uk and CURIA public source pages for GDPR Article 15 access-right materials. The layer contains 19 public records across three files.

The source-bound EU public-source packet qualified as normative material screening because it included the relevant high-authority materials, limiting provisions, source locators and reviewable output-source links. The U.S. and English issue-search outputs were capped at reference information because the visible returned records omitted high-authority or counter-material records for the defined issue. The procedural-status result is sharper than a relevance score: public legal records can be source-reconstructable yet fail normative-screening qualification when the output does not preserve the issue-specific authority and counter-material chain.

\subsection{Public retrieval benchmark}

To make that result less dependent on isolated examples, the harness also freezes 30 public retrieval outputs and scores them against endpoint-compatible authority sets. The benchmark covers five issue families: U.S. agency deference after \textit{Loper Bright}, English mesothelioma causation after \textit{Fairchild}, Canadian administrative-law standard of review after \textit{Vavilov}, German/EU right-to-be-forgotten review and EU GDPR Article 15 access rights. It uses CourtListener, the National Archives, Supreme Court of Canada Lexum search, OpenLegalData and CURIA public search or list endpoints. The benchmark contains 225 returned records. It is not a hand-built positive control: the returned top-$k$ records are whatever the public endpoints exposed for the committed queries. Its gold sets are matched to endpoint scope, so case-law endpoints are not penalized for failing to return statutes; mixed statute/case screening is tested in the source-bound issue packets.

All 30 outputs were capped at reference information. Mean high-authority recall was 0.26 and mean counter-material recall was 0.00. Source-reconstructable public search output is therefore not the same as normative material screening. A system can expose real cases, stable URLs and rank order while still failing the procedural requirement that high-authority and limiting materials for a defined legal issue be made visible together.

\subsection{Raw model-output pilot}

The harness was also run on ten raw legal-authority recommendations generated by a single upstream model configuration, Codex GPT-5.5 with extra-high reasoning. The prompts covered the same three issue families as the positive-control packets: U.S. agency deference after \textit{Loper Bright}, English mesothelioma causation after \textit{Fairchild}, and EU GDPR Article 15 access rights. The model was not allowed to browse or use tools. This isolates the status question: whether fluent model text with strong authority coverage can itself claim normative screening status.

It cannot. In the coded raw-output packets, all ten outputs identified the expected high-authority and limiting materials in the issue families, producing mean upstream recall of 1.00 and counter-authority recall of 1.00. Yet every claimed normative-screening output was downgraded to reference information because the citations were marked as requiring verification rather than bound to official sources, corpus manifests, pinpoint locators and a contestation pathway. Their source-tag coverage was 1.00, but procedural-source coverage was 0.00.

The repair intervention then applies the opposite condition to the same authority sets. A validator checks each output-source link against an issue manifest, locator, hashed source-support excerpt, contradiction pattern, procedural source tag and issue-set coverage before assigning repair scores; when a limiting material is already retrieved but outside the review window, a separate rank-salience repair moves that material into view. A transcript verifier first checks that all 50 scenario locators appear in the ten committed raw-output sections. A separate public source-text layer checks the manifest support terms against extracted source snapshots; all 30 support items were externally anchored across 30 records with text snapshots. All ten repaired variants qualified as normative material screening and none became decision-support reasons. This paired result is the paper's central rule in executable form: authority coverage is not procedural status, but validator-backed source support and visible counter-material can make authority coverage procedurally screenable when contestability and logging gates are satisfied.

The evidence-ladder intervention then varies seven conditions over ten model-output variants: raw unverified output, source-bound output without counter-material, source-bound output without contestability, source-bound output without logging, contestable screening, authorized decision-support adoption and unauthorized external action. The ladder produced the predicted status pattern: 30 reference-information outputs, 10 professional-support outputs, 10 normative-screening outputs, 10 decision-support reasons and 10 outputs with no external legal effect. This is the controlled result: generated legal content changes procedural status only when the legal-material chain, rank salience or adoption gate changes.

The same ten outputs were also converted into sixty adversarial repair variants. The attacks introduced locator mismatches, unsupported claims, contradiction-pattern claims, out-of-manifest sources, missing output-source links and counter-material omissions. All sixty adversarial variants were rejected despite mean upstream recall of 1.00: forty were capped at reference information, while twenty unsupported or contradictory source-support variants were withdrawn from external legal effect. The blocked-failure distribution across the full matrix shows why high-performing outputs failed: source-attribution gaps dominated, followed by summary distortion, counter-material suppression, authority omission and smaller numbers of contestability and invalid-authority failures.

\begin{table}[htbp]
\centering
\caption{Raw model-output pilot}
\small
\renewcommand{\arraystretch}{1.1}
\begin{tabularx}{\textwidth}{>{\bfseries}p{2.8cm}p{1.2cm}p{1.6cm}p{1.7cm}p{1.0cm}Y}
\toprule
Issue family & N & Mean score & Upstream recall & PSC & Allowed status \\
\midrule
U.S. agency deference after \textit{Loper Bright} & 3 & 9 & 1.00 & 0.00 & Reference information \\
English mesothelioma causation after \textit{Fairchild} & 3 & 9 & 1.00 & 0.00 & Reference information \\
EU GDPR Article 15 access rights & 4 & 9 & 1.00 & 0.00 & Reference information \\
\bottomrule
\end{tabularx}
\end{table}

To test mixed-authority normative material screening itself, the harness includes five public source-screening packets: U.S. administrative law after \textit{Loper Bright}, English-law mesothelioma causation after \textit{Fairchild}, EU GDPR Article 15 access-right interpretation, Canadian administrative-law standard of review after \textit{Vavilov}, and German/EU fundamental-rights review after the \textit{Right to be Forgotten} decisions. Each packet defines high-authority materials, limiting or counter-materials, invalidated or historical treatments, source tags and output-to-source links at the level of individual output units. These are mixed statute/case/source construct tests, not a benchmark leaderboard: they ask whether a source-bound, issue-defined authority packet can obtain screening status under the protocol. All five qualified as normative material screening, not as decision-support reasons, because no authorized decision-maker adopted them as reasoned outcomes.

\begin{table}[htbp]
\centering
\caption{Mixed-authority public source-screening packets}
\scriptsize
\renewcommand{\arraystretch}{1.1}
\begin{tabularx}{\textwidth}{>{\bfseries}Yp{1.8cm}p{2.3cm}p{2.0cm}p{1.8cm}}
\toprule
Issue packet & Profile & Authority set & Metrics & Status \\
\midrule
U.S. agency deference after \textit{Loper Bright} & Common law & H=3; C=2; invalid=1 treatment & Auth.=1.00; CAR=1.00; EF=1.00 & Normative screening \\
English mesothelioma causation after \textit{Fairchild} & Common law & H=3; C=2; invalid=1 treatment & Auth.=1.00; CAR=1.00; EF=1.00 & Normative screening \\
EU GDPR Article 15 access rights & Civil law/EU & H=3; C=2; invalid=2 treatments & Auth.=1.00; CAR=1.00; EF=1.00 & Normative screening \\
Canada standard of review after \textit{Vavilov} & Common law & H=3; C=2; historical=1 treatment & Auth.=1.00; CAR=1.00; EF=1.00 & Normative screening \\
Germany/EU right-to-be-forgotten review & Civil law/EU & H=3; C=2; cross-system=2 treatments & Auth.=1.00; CAR=1.00; EF=1.00 & Normative screening \\
\bottomrule
\end{tabularx}
\end{table}

The ablation suite then removes one procedural condition at a time from those packets. Twenty ablations were generated: high-authority omission, counter-material suppression, unverified source tags and missing authorized-adoption gates for each issue family. Fifteen were capped at reference information through downgrade or suspension. The remaining five claimed decision-support status but were capped at normative screening because the packet lacked authorized institutional adoption. This is the crucial negative test: a packet that is otherwise strong on source links and authority coverage still loses status when the missing condition is procedural rather than semantic.

The metric-separation layer uses the 185 packets with complete upstream precision, recall and F1 scores to test whether conventional retrieval metrics predict procedural qualification. They do not. The point-biserial correlation between upstream recall and procedural qualification was 0.05. A recall threshold of 0.8 identified every procedurally qualified packet but also produced 136 false positives, yielding precision 0.25. The statistical robustness layer then ran 1,000 bootstrap resamples and 1,000 label permutations. The 95\% bootstrap interval for recall-threshold precision was 0.19--0.32; the interval for high-recall blocked rate was 0.68--0.81; full-gate precision and specificity stayed at 1.00; and the recall/status permutation test returned $p=0.538$. The audit gate cascade then removes the high-recall false positives by adding score eligibility, source binding, review-role readiness and absence of failure caps. This result is the empirical counterpart to the formal metric non-equivalence invariant: retrieval success may be necessary in many settings, but it is not a status-conferring condition.

The gate-ablation layer then tests gate necessity across the outputs that did qualify. For each of the 46 packets that reached normative screening or decision-support status, the harness generated counterfactual variants removing the evidence packet, replacing procedural source tags with verification-needed tags, omitting high-authority materials, omitting counter-materials, removing the review gate or imposing a professional-support role cap. For the 12 decision-support packets, it also removed the decision-adoption reason. All 288 ablations fell below the corresponding procedural level. The result matters because it shows that qualification is not carried by an isolated issue packet or a high total score; it depends on the procedural chain as a whole.

The repair-frontier layer applies the converse test to outputs that were blocked below their claimed high-status level. For each blocked normative-screening or decision-support claim, the harness enumerates counterfactual combinations of seven artifact-gate families: audit-vector completion, authority-packet completion, rank-salience completion, source-binding completion, review-contestability completion, role-cap completion and failure-flag resolution; decision-support claims also include adoption completion. The layer evaluated 4,418 counterfactual variants across 176 blocked high-status claims. All 176 were repairable by the declared gate families. The minimal-frontier distribution was concentrated rather than arbitrary: 100 claims required one gate family, 34 required two, 39 required three and three required four. Source-binding completion appeared most often in minimal frontiers, followed by authority-packet completion and rank-salience completion. The practical result is that downgrade is not a dead end: the protocol identifies the smallest class of missing artefact that would have to be supplied before a blocked output could be reconsidered.

The jurisdiction-profile layer targets the objection that a common audit vector may hide common-law assumptions. The harness checked 217 high-status claims and found 217/217 profile checks supported. It then mutated all 46 qualified packets: replacing specific assumptions with the valid generic profile preserved status in 46/46 cases, removing the assumption downgraded 46/46 cases and substituting a mismatched profile downgraded 46/46 cases, for 138/138 profile mutations passed. The result is not jurisdiction neutrality; it is profile-bound portability.

The ranking-visibility layer targets the paper's visibility claim without treating salience as a hidden merits score. Across 214 high-status claims, 820 rank-window visibility checks produced a counter-material visibility curve: 14/214 at rank 1, 70/207 by rank 2, top-3 counter-material visibility was 167/201, and 173/198 by rank 4. The median first counter-material rank was 3.0 and the mean reciprocal first-counter rank was 0.43. In the 70 packets with enough non-counter material to run a true rank-order intervention, the counterfactual moved counter-material below the visibility window while preserving authority coverage, counter-material recall, evidence fidelity and procedural source-tag coverage. All 70 were downgraded below normative screening, and drifted top-3 counter-material visibility was 0/70. The result is narrower than a field study but sharper than set recall: the same authority set can be present while rank salience remains independently status-relevant.

The status-certificate layer turns each allocation into a replayable audit object. For every packet, the harness records the scenario hash, policy thresholds, claimed status, jurisdiction profile, score vector, score-based candidate status, role cap, failure cap, disposition, final status and metric bundle. Replaying those certificates verified 230/230 status certificates and 2,990/2,990 field checks. This makes the status assignment inspectable without rerunning the whole experimental narrative.

The final layers address coding uncertainty. First, all 230 committed scenario packets were re-scored under strict and lenient policies. Status allocation was stable for 218 of 230 packets; weighted status agreement was 1.00 for base versus strict and 0.98 for base versus lenient. Second, the harness perturbed all six audit scores 57,500 times under a fixed-seed annotation-noise model. Sample-level status stability was 0.936; among the 46 qualified packets, exact-status stability was 0.876 and high-status stability was 0.928. The instability is informative rather than hidden: the boundary cases were concentrated around professional-support/screening and screening/decision-support gates, where one-point coding changes affect the threshold but do not defeat the non-compensatory gate logic. Third, the same 230 packets were converted into score-blinded packet files that removed original scores, expected outcomes, source paths and manual failure flags. Two codebook-based coding passes then applied fresh scores to the packets. This is a codebook reproducibility check, not an external annotation study: treatment names and packet facts remain visible because coders must still read the legal-output evidence. The passes reached 0.99 exact and weighted inter-coder status agreement; against the base harness allocation, the weaker coder reached 0.95 exact agreement and 0.97 weighted agreement. The disagreements identify the fragile gates: decision-support adoption, deliberately thin stress packets and ablations near the screening threshold.

\section{Implications and Limitations}

The framework has three implications for legal AI evaluation. First, model performance is not procedural permission. A system may meet technical retrieval benchmarks while failing procedural qualification because the output cannot be reconstructed, source-tagged or contested. Second, human oversight is not magic. Oversight must leave review status, reasons, rejection records and adoption trails; a human signature without verifiable review should not upgrade an output's status. Third, accountability should follow control. Developers, deployers, professional users and institutional adopters control different risks and should carry corresponding explanation duties.

The principal limitation is institutional realism rather than executable coverage. The current validation combines public search outputs, model-output interventions, transcript anchors, public source-text anchors, adversarial repair attacks, metric-separation tests, gate-ablation tests, repair-frontier diagnosis, jurisdiction-profile mutation, ranking-visibility counterfactuals, status-certificate replay, scenario tests, artifact checks, score-noise perturbation and score-blinded codebook coding rather than a deployed-system field study with external human annotators. It shows how legal-output status can be audited above heterogeneous upstream systems and how robust the status rules are to source-binding attacks, missing-gate repairs, profile mismatch, ranking drift and coding uncertainty, while leaving adoption effects, user behavior and larger independent annotation for later work.

The next empirical step is therefore clear rather than open-ended. Deployed case recommendation systems and larger legal-output benchmarks should be scored for corpus verifiability, counter-material presentation, ranking traceability, adoption logging and contestation support, and those scores should be compared with retrieval metrics and user behavior. The question for future work is not only whether legal AI systems are accurate, but whether their outputs are legally usable.

\section{Conclusion}

Status follows auditability, not intelligence. That is the paper's central rule. A case recommendation system may be technically sophisticated, fluent or empirically strong and still have no procedural status if its legal-material chain cannot be reconstructed and challenged. Conversely, an output does not need to be a legal source or autonomous decision-maker to matter procedurally. Once it structures the visibility and ranking of authorities, it occupies the middle category of normative material screening.

The proposed framework makes that category operational. It separates output effect from system role and links both to source, corpus, ranking, counter-material, contestability, adoption and audit requirements. It allows an output to remain internal reference information, become professional support, qualify as normative material screening, or enter a decision reason only through accountable human adoption. It also supplies downgrade and withdrawal rules when the evidentiary and procedural chain fails.

The contribution is therefore not another general claim that AI cannot be a judge or a legal source. It is a stricter evaluation logic for legal AI outputs: no source chain, no status; no counter-material pathway, no screening qualification; no adoption record, no decision reason.

\section*{Data and Code Availability}

The companion artifact contains the audit harness, scenario packets, public snapshots, raw model-output transcript, validation scripts and generated Markdown/JSON reports. The primary replication command is \path{python scripts/run_full_validation.py}; it reruns the full validation matrix and writes the aggregate report under \path{experiments/full_validation/results/}. The artifact is designed for deterministic review without LLM API access because public-source snapshots and raw model outputs are committed. Fresh public collection is available only through explicit refresh commands and is not required to reproduce the reported results.

\begin{thebibliography}{99}

\bibitem[Ashley(1992)]{ashley1992}
Ashley KD (1992) Case-based reasoning and its implications for legal expert systems. \textit{Artificial Intelligence and Law} 1(2):113--208. \url{https://doi.org/10.1007/BF00114920}

\bibitem[Bench-Capon and Sartor(2003)]{benchcapon2003}
Bench-Capon TJM, Sartor G (2003) A model of legal reasoning with cases incorporating theories and values. \textit{Artificial Intelligence} 150(1--2):97--143. \url{https://doi.org/10.1016/S0004-3702(03)00108-5}

\bibitem[Binns et~al.(2018)]{binns2018}
Binns R, Van Kleek M, Veale M, Lyngs U, Zhao J, Shadbolt N (2018) ``It's reducing a human being to a percentage'': Perceptions of justice in algorithmic decisions. In: \textit{Proceedings of the 2018 CHI Conference on Human Factors in Computing Systems}. ACM, pp 1--14. \url{https://doi.org/10.1145/3173574.3173951}

\bibitem[Chalkidis et~al.(2022)]{chalkidis2022}
Chalkidis I, Jana A, Hartung D, Bommarito M, Androutsopoulos I, Katz DM, Aletras N (2022) LexGLUE: A benchmark dataset for legal language understanding in English. In: \textit{Proceedings of the 60th Annual Meeting of the Association for Computational Linguistics}. ACL, pp 4310--4330. \url{https://doi.org/10.18653/v1/2022.acl-long.297}

\bibitem[Citron and Pasquale(2014)]{citron2014}
Citron DK, Pasquale F (2014) The scored society: Due process for automated predictions. \textit{Washington Law Review} 89(1):1--33. \url{https://digitalcommons.law.uw.edu/wlr/vol89/iss1/2/}

\bibitem[Dahl et~al.(2024)]{dahl2024}
Dahl M, Magesh V, Suzgun M, Ho DE (2024) Large legal fictions: Profiling legal hallucinations in large language models. \textit{Journal of Legal Analysis} 16(1):64--93. \url{https://doi.org/10.1093/jla/laae003}

\bibitem[Dung(1995)]{dung1995}
Dung PM (1995) On the acceptability of arguments and its fundamental role in nonmonotonic reasoning, logic programming and n-person games. \textit{Artificial Intelligence} 77(2):321--357. \url{https://doi.org/10.1016/0004-3702(94)00041-X}

\bibitem[Gebru et~al.(2021)]{gebru2021}
Gebru T, Morgenstern J, Vecchione B, Vaughan JW, Wallach H, Daum\'{e} H III, Crawford K (2021) Datasheets for datasets. \textit{Communications of the ACM} 64(12):86--92. \url{https://doi.org/10.1145/3458723}

\bibitem[Goebel et~al.(2024)]{goebel2024}
Goebel R, Kano Y, Kim MY, Rabelo J, Satoh K, Yoshioka M (2024) Overview and discussion of the Competition on Legal Information, Extraction/Entailment (COLIEE) 2023. \textit{Review of Socionetwork Strategies} 18:27--47. \url{https://doi.org/10.1007/s12626-023-00152-0}

\bibitem[Gordon et~al.(2007)]{gordon2007}
Gordon TF, Prakken H, Walton D (2007) The Carneades model of argument and burden of proof. \textit{Artificial Intelligence} 171(10--15):875--896. \url{https://doi.org/10.1016/j.artint.2007.04.010}

\bibitem[Guha et~al.(2023)]{guha2023}
Guha N, Nyarko J, Ho DE, Re C, Chilton A, Narayana A, Chohlas-Wood A, Peters A, Waldon B et al (2023) LegalBench: A collaboratively built benchmark for measuring legal reasoning in large language models. In: \textit{Advances in Neural Information Processing Systems 36}. \url{https://papers.nips.cc/paper_files/paper/2023/hash/89e44582fd28ddfea1ea4dcb0ebbf4b0-Abstract-Datasets_and_Benchmarks.html}

\bibitem[Hou et~al.(2025)]{hou2025}
Hou AB, Weller O, Qin G, Yang E, Lawrie D, Holzenberger N, Blair-Stanek A, Van Durme B (2025) CLERC: A dataset for U.S. legal case retrieval and retrieval-augmented analysis generation. In: \textit{Findings of the Association for Computational Linguistics: NAACL 2025}. ACL, pp 7913--7928. \url{https://doi.org/10.18653/v1/2025.findings-naacl.441}

\bibitem[Kroll et~al.(2017)]{kroll2017}
Kroll JA, Huey J, Barocas S, Felten EW, Reidenberg JR, Robinson DG, Yu H (2017) Accountable algorithms. \textit{University of Pennsylvania Law Review} 165(3):633--705. \url{https://scholarship.law.upenn.edu/penn_law_review/vol165/iss3/3/}

\bibitem[Magesh et~al.(2025)]{magesh2025}
Magesh V, Surani F, Dahl M, Suzgun M, Manning CD, Ho DE (2025) Hallucination-free? Assessing the reliability of leading AI legal research tools. \textit{Journal of Empirical Legal Studies} 22(2):216--242. \url{https://doi.org/10.1111/jels.12413}

\bibitem[Mitchell et~al.(2019)]{mitchell2019}
Mitchell M, Wu S, Zaldivar A, Barnes P, Vasserman L, Hutchinson B, Spitzer E, Raji ID et al (2019) Model cards for model reporting. In: \textit{Proceedings of the Conference on Fairness, Accountability, and Transparency}. ACM, pp 220--229. \url{https://doi.org/10.1145/3287560.3287596}

\bibitem[Mökander(2023)]{mokander2023}
Mökander J (2023) Auditing of AI: Legal, ethical and technical approaches. \textit{Digital Society} 2:49. \url{https://doi.org/10.1007/s44206-023-00074-y}

\bibitem[Prakken and Sartor(2015)]{prakken2015}
Prakken H, Sartor G (2015) Law and logic: A review from an argumentation perspective. \textit{Artificial Intelligence} 227:214--245. \url{https://doi.org/10.1016/j.artint.2015.06.005}

\bibitem[Raji et~al.(2020)]{raji2020}
Raji ID, Smart A, White RN, Mitchell M, Gebru T, Hutchinson B, Smith-Loud J, Theron D, Barnes P (2020) Closing the AI accountability gap: Defining an end-to-end framework for internal algorithmic auditing. In: \textit{Proceedings of the 2020 Conference on Fairness, Accountability, and Transparency}. ACM, pp 33--44. \url{https://doi.org/10.1145/3351095.3372873}

\bibitem[Rissland and Daniels(1995)]{rissland1995}
Rissland EL, Daniels JJ (1995) A hybrid CBR-IR approach to legal information retrieval. In: \textit{Proceedings of the 5th International Conference on Artificial Intelligence and Law}. ACM, pp 52--61. \url{https://doi.org/10.1145/222092.222125}

\bibitem[Shao et~al.(2020)]{shao2020}
Shao Y, Mao J, Liu Y, Ma W, Satoh K, Zhang M, Ma S (2020) BERT-PLI: Modeling paragraph-level interactions for legal case retrieval. In: \textit{Proceedings of the Twenty-Ninth International Joint Conference on Artificial Intelligence}. IJCAI, pp 3501--3507. \url{https://doi.org/10.24963/ijcai.2020/484}

\bibitem[Turtle(1995)]{turtle1995}
Turtle H (1995) Text retrieval in the legal world. \textit{Artificial Intelligence and Law} 3:5--54. \url{https://doi.org/10.1007/BF00877694}

\bibitem[van Opijnen and Santos(2017)]{vanopijnen2017}
van Opijnen M, Santos C (2017) On the concept of relevance in legal information retrieval. \textit{Artificial Intelligence and Law} 25:65--87. \url{https://doi.org/10.1007/s10506-017-9195-8}

\bibitem[Veale et~al.(2018)]{veale2018}
Veale M, Van Kleek M, Binns R (2018) Fairness and accountability design needs for algorithmic support in high-stakes public sector decision-making. In: \textit{Proceedings of the 2018 CHI Conference on Human Factors in Computing Systems}. ACM, pp 1--14. \url{https://doi.org/10.1145/3173574.3174014}

\bibitem[Wachter et~al.(2017)]{wachter2017}
Wachter S, Mittelstadt B, Floridi L (2017) Why a right to explanation of automated decision-making does not exist in the General Data Protection Regulation. \textit{International Data Privacy Law} 7(2):76--99. \url{https://doi.org/10.1093/idpl/ipx005}

\bibitem[Zheng et~al.(2021)]{zheng2021}
Zheng L, Guha N, Anderson BR, Henderson P, Ho DE (2021) When does pretraining help? Assessing self-supervised learning for law and the CaseHOLD dataset of 53,000+ legal holdings. In: \textit{Proceedings of the 18th International Conference on Artificial Intelligence and Law}. ACM, pp 159--168. \url{https://doi.org/10.1145/3462757.3466088}

\bibitem[Zheng et~al.(2025)]{zheng2025}
Zheng L, Guha N, Arifov J, Zhang S, Skreta M, Manning CD, Henderson P, Ho DE (2025) A reasoning-focused legal retrieval benchmark. In: \textit{Proceedings of the 4th ACM Symposium on Computer Science and Law}. ACM, pp 169--193. \url{https://doi.org/10.1145/3709025.3712219}

\end{thebibliography}

\end{document}
